\begin{chapintro}
\kkref{sec:04-2} で作った 10 試験の 2×2 の数を、一つの効果の推定値にまとめます。分散、対数、オッズといった計算の前提を精度管理の言葉で確かめたあと、効果をどの指標で表すか、試験にどう重みを付けるか、固定効果モデルとランダム効果モデルのどちらを使うかを決め、R の meta パッケージで計算します。結果はフォレストプロットで読みます。τ² の推定法の詳細と L'Abbé プロットは「発展」の印をつけた部分で扱うので、初めて読むときは飛ばして構いません。
\end{chapintro}

\section{計算の前に知っておくこと}
\label{sec:05-1}

この章の計算は、分散、標準誤差、対数、オッズという四つの言葉の上に成り立っています。どれも精度管理で使う感覚に置き換えられます。

\subsection{ばらつきを表す量}
\label{sec:05-1-1}

標準偏差（SD）は、測定値のばらつきの大きさです。同じ管理試料を測ると、値は測るたびに SD ほどぶれます。分散は SD を 2 乗した値です。標準誤差は、個々の測定値ではなく、平均やリスク比のような推定値そのものの不確かさを表します。管理試料を 1 回だけ測った値より、20 回測った平均のほうがぶれが小さいのと同じで、標準誤差は人数や測定回数が多いほど小さくなります。平均の標準誤差は SD を測定回数の平方根で割った値なので、20 回の平均のぶれは 1 回の値の約 4.5 分の 1 です。試験でも、人数が多いほど推定値のぶれは小さくなります。

\subsection{対数と指数}
\label{sec:05-1-2}

リスク比 0.5 と 2 は、向きが逆で大きさの同じ効果です。ところが数直線の上では、1 からの距離が 0.5 と 1 で違います。対数をとると、0.5 は \mbox{−0.69}、1 は 0、2 は 0.69 になり、0 を中心に左右対称に並びます（表 5.1）。log は掛け算を足し算に変える計算で、log(1) = 0 です。対数をとった値は、exp（指数）をとると元の比に戻ります。

\begin{kkcenter}
\kkfit{\begin{tabular}{lrrrrr}
\toprule
リスク比 & 0.25 & 0.5 & 1 & 2 & 4 \\
1 からの距離 & 0.75 & 0.5 & 0 & 1 & 3 \\
対数をとった値 & −1.39 & −0.69 & 0 & 0.69 & 1.39 \\
\bottomrule
\end{tabular}}
\kkcap{表 5.1　リスク比と、その対数をとった値}
\end{kkcenter}

\subsection{重みを分散の逆数にする理由}
\label{sec:05-1-3}

精度管理では、CV の小さい、ぶれの少ない測定ほど信用します。メタ解析も同じで、ぶれの小さい試験を重く見ます。ぶれの大きさを表すのが分散なので、分散の逆数を重みにします。降圧薬 A の Endo 2015（各群 90 人）と Garcia 2018（各群 1600 人）では、リスク比の対数の分散が約 0.23 と約 0.012 で、Garcia の重みは Endo の約 18 倍になります。

\subsection{オッズ}
\label{sec:05-1-4}

オッズは、ある出来事が起きる割合 p を、起きない割合 1 − p で割った p/(1 − p) です。発症割合が 0.2 なら、オッズは 0.2/0.8 = 0.25 で、発症する人 1 人に対して発症しない人が 4 人いることを表します。オッズは \kkref{sec:09-4} と \kkref{sec:10-3} でも使います。

\section{効果を表す指標}
\label{sec:05-2}

メタ解析で統合するのは、試験ごとの効果の大きさです。降圧薬 A の試験では、アウトカムが脳卒中を起こしたか起こさなかったかの 2 値なので、二つの群の発症割合を比べる指標を使います。代表的なものが三つあります。

\begin{description}
\item[リスク比（RR）] 介入群の発症割合を対照群の発症割合で割った値です。架空の Garcia 2018 では降圧薬 A 群が 1600 人中 121 人（7.6\%）、プラセボ群が 1600 人中 186 人（11.6\%）なので、リスク比は 7.6 ÷ 11.6 で 0.65 になります。1 より小さければ介入群で発症が少なかったことを表し、0.65 は発症が 35\% 少ないと読めます。
\item[オッズ比（OR）] オッズ（\kkref{sec:05-1-4}）を群どうしで割った値です。発症がまれなアウトカムではリスク比とほぼ同じ値になりますが、発症割合が高くなるほどリスク比より 1 から遠い値になります。臨床の言葉に直しにくいので、RCT のメタ解析ではリスク比を選ぶことが多くなります。
\item[リスク差（RD）] 発症割合の差です。Garcia 2018 では 121/1600 − 186/1600 で約 −4.1 ポイント（7.56\% − 11.63\%）になります。1000 人あたり何人減るかという形で判断に結びつけやすい指標ですが、対照群の発症割合が試験ごとに違うと試験間で大きくばらつくので、統合にはあまり使いません。
\end{description}

アウトカムが血圧の値や検査値のような連続量なら、平均の差（平均差、MD）を使います。生活の質を別々の質問票で測った試験のように、試験ごとに尺度の単位が違うときは、平均差を標準偏差で割った標準化平均差（SMD）を使います。

リスク比とオッズ比は、そのままではなく対数をとってから統合します（対数は \kkref{sec:05-1-2} で説明しました）。比は 1 を中心に左右対称ではないからです。統合した結果は最後に指数をとって比に戻します。フォレストプロットの横軸が対数目盛になっているのはこのためです。

この冊子ではリスク比を使います。降圧薬 A を使うかどうかを考える人にとって、何割減るかが最も読み取りやすく、対照群の発症割合が試験ごとに約 5〜13\% と違っていても比はばらつきにくいからです。

\section{重み付き平均という考え方}
\label{sec:05-3}

10 試験のリスク比を単純に平均すると、各群 90 人の Endo 2015 と各群 3000 人の Ivanov 2021 を同じ重さで扱うことになります。大きな試験ほど偶然のぶれが小さく、推定値を信用できるので、重みを大きくするのが自然です。重みの考え方が分かりにくいときは、先に \kkref{sec:05-1-3} を読み直してください。

メタ解析の標準的な重みは、推定値の分散の逆数です。この重みの付け方を逆分散法と呼びます。

対数リスク比の分散は、主に事象の数で決まります。人数が多くても脳卒中がほとんど起きなければ、推定値は安定しません。Garcia 2018 は両群で 307 件の脳卒中が起きたので分散が小さく、固定効果モデルでは 10 試験の中で最も重い 23.9\% の重みを受けます。Endo 2015 は 18 件しかないので、重みは 1.3\% にとどまります。統合した値は、この重みで対数リスク比を平均したものです。

\begin{KKfbox9ao}[計算の補足]
対数リスク比の分散は、1/（介入群の事象数）− 1/（介入群の人数）+ 1/（対照群の事象数）− 1/（対照群の人数）で計算できます。Endo 2015 なら 1/6 − 1/90 + 1/12 − 1/90 を計算し、Garcia 2018 なら同じ式に 121、1600、186、1600 を入れます。標準誤差は分散の平方根で、Endo 2015 は約 0.48、Garcia 2018 は約 0.11 です。
\end{KKfbox9ao}

\section{固定効果モデルとランダム効果モデル}
\label{sec:05-4}

重みを付けて平均するとき、試験の背後にある真の効果をどう考えるかで、二つのモデルに分かれます。

同じ管理試料を複数の施設で測る場面で考えます。固定効果モデルは、どの施設も同じ真の値を測っていて、施設ごとの結果の違いは測定の偶然誤差だけだと考えます。試験でいえば、すべての試験が同じ一つの真の効果を推定していて、試験ごとの結果が違うのは偶然のぶれだけだという前提です。R の meta パッケージは、このモデルを common effect model（共通効果モデル）と表示します。

ランダム効果モデルは、施設ごとに校正や試薬ロットの違いによる系統的なずれがあり、施設が測っている値そのものが少しずつ違うと考えます。試験でも、対象の年齢、併用している治療、追跡期間が違えば、真の効果が少しずつ違います。このモデルでは、試験ごとの真の効果が一つの分布から出てきたとみなし、その分布の平均を推定します。真の効果の散らばりの大きさを研究間分散と呼び、τ²（タウ二乗）で表します。τ² は施設間のずれの大きさに当たり、各施設の中の偶然誤差、つまり各試験の分散とは別のものです。

\begin{kkcenter}
\begin{kkfitpic}
\begin{tikzpicture}[x=1mm,y=1mm, font=\footnotesize]
  \node[font=\small\gtfamily, anchor=west] at (0,48) {固定効果モデル};
  \draw[kkdark, line width=.6pt] (0,8) -- (52,8);
  \draw[kkwarm, line width=1pt] (26,8) -- (26,40);
  \node[kkwarm, anchor=south] at (26,40) {真の効果は一つ};
  \foreach \x/\y in {20/34, 29/28, 24/22, 31/16, 22/12}{
    \draw[kkaccent, line width=.6pt] (\x-4,\y) -- (\x+4,\y);
    \fill[kkaccent] (\x,\y) circle (1.1);
  }
  \node[anchor=north, align=center] at (26,6) {試験ごとの差は\\偶然のぶれだけ};
  \begin{scope}[xshift=62mm]
  \node[font=\small\gtfamily, anchor=west] at (0,48) {ランダム効果モデル};
  \draw[kkdark, line width=.6pt] (0,8) -- (52,8);
  \draw[kksecondary, fill=kkpale, line width=.6pt]
    plot[domain=8:44, samples=60] (\x, {8 + 26*exp(-((\x-26)^2)/(2*6^2))});
  \draw[kkwarm, line width=1pt, dashed] (26,8) -- (26,40);
  \node[kkwarm, anchor=south] at (26,40) {真の効果の平均};
  \foreach \t/\o/\y in {18/15/31, 30/33/25, 23/21/19, 34/37/14}{
    \fill[kkdark] (\t,8) circle (.8);
    \draw[kkdark, densely dotted] (\t,8) -- (\o,\y);
    \draw[kkaccent, line width=.6pt] (\o-4,\y) -- (\o+4,\y);
    \fill[kkaccent] (\o,\y) circle (1.1);
  }
  \node[anchor=north, align=center] at (26,6) {試験ごとの真の効果（黒点）が\\違い、そこに偶然のぶれが加わる};
  \end{scope}
\end{tikzpicture}
\end{kkfitpic}
\kkcap{図 5.1　二つのモデルの考え方。青い点と横線は試験ごとの推定値と信頼区間}
\end{kkcenter}

二つのモデルの違いは、重みの付け方に表れます。ランダム効果モデルでは、各試験の分散に τ² を足したものの逆数を重みにします。τ² はどの試験にも同じだけ足されるので、もともと分散が小さかった大きな試験ほど重みが大きく目減りし、試験間の重みの差が縮まります。Endo 2015 の分散 0.23 に τ² の 0.034 を足すと 0.26、Garcia 2018 の 0.012 に足すと 0.046 で、重みの比は約 18 倍から約 6 倍に縮みます。重みの割合でみると、Garcia 2018 の重みが 23.9\% から 16.5\% に下がり、Endo 2015 は 1.3\% から 2.9\% に上がります。小さな試験の重みが相対的に増えるので、\kkref{sec:07-1} で扱う小規模研究効果があるときには、この性質が結果を動かします。

\kkfig[.85\linewidth]{x1_weights.pdf}{図 5.2　固定効果モデルとランダム効果モデルでの重み。各試験の上の棒が固定効果、下の棒がランダム効果。試験は合計人数の少ない順に下から並べた}

図 5.2 の横軸は、10 試験の合計を 100 としたときの各試験の重み（\%）です。各試験に 2 本の棒があり、薄い青が固定効果モデル、濃い青がランダム効果モデルの重みを表します。下の 3 試験（Endo 2015、Aoki 2009、Hayashi 2019）はどれも合計 600 人以下の小さな試験で、濃い青の棒が薄い青より長くなっています。上から 3 本の大きな試験（Ivanov 2021、Dubois 2014、Garcia 2018）では逆に、濃い青の棒が短くなっています。

もう一つ読み取れることは、棒の長さが人数の順になっていないことです。Garcia 2018 は合計 3,200 人で、Dubois 2014（5,000 人）や Ivanov 2021（6,000 人）より少ないのに、固定効果モデルでは最も重くなっています。脳卒中の件数が多く（307 件）、推定が安定しているからです。重みは人数ではなく、推定のぶれの小ささで決まります。ランダム効果モデルで重みの差が縮んでも、並び順は大きくは変わりません。

どちらを主解析にするかは、データを見る前に計画の段階で決めます。臨床試験は対象や介入の条件が試験ごとに違うのが普通なので、この冊子ではランダム効果モデルを主解析にします。ランダム効果モデルが推定するのは真の効果の平均で、どの集団でもその効果が得られるという意味ではありません（予測区間は \kkref{sec:06-3} で説明します）。

\section{R で計算する準備}
\label{sec:05-5}

ここから R を使います。R を初めて使う読者は、付録 A の手順で R と RStudio を入れ、作業フォルダ（中に data、fig、out の 3 つのフォルダ）を作ってから読み進めてください。

RStudio の File メニューから New File、R Script と選ぶとスクリプトの画面が開きます。命令を書いた行で Command と Return（Windows では Ctrl と Enter）を押すと、その行が実行され、結果が Console に出ます。\# から右は R が読み飛ばす注釈なので、写すときに省いてもかまいません。

最初に、メタ解析のパッケージを読み込みます。インストール（付録 A）は一度だけ、読み込みは R を起動するたびに行います。

\begin{rcode}
library(meta)      # メタ解析の関数をまとめたパッケージ
\end{rcode}

次に、第 4 章で作った CSV ファイルを読み込みます。read.csv は CSV を表の形のデータ（データフレーム）として読み込む関数です。左向きの矢印 \texttt{<-} は、右側の結果を左側の名前で保存するという意味で、この例では読み込んだ表を bp という名前にしています。

\begin{rcode}
bp <- read.csv("data/bp_stroke.csv",
               fileEncoding = "UTF-8")
str(bp)            # 列の名前と型を確かめる
\end{rcode}

\begin{rout}
'data.frame':	10 obs. of  7 variables:
 $ study: chr  "Aoki 2009" "Berger 2011" "Chen 2012" ...
 $ pop  : chr  "高齢者" "一般" "高齢者" "一般" ...
 $ rob  : chr  "懸念あり" "低" "懸念あり" "低" ...
 $ ev_t : int  7 48 44 134 6 32 121 10 112 43
 $ n_t  : int  150 1200 420 2500 90 800 1600 300 3000 650
 $ ev_c : int  17 65 39 145 12 38 186 23 139 83
 $ n_c  : int  150 1200 420 2500 90 800 1600 300 3000 650
\end{rout}

10 行 7 列の表として読み込めました。ev\_t と n\_t が介入群の事象数と人数、ev\_c と n\_c が対照群の事象数と人数です。ファイルが見つからないときは作業フォルダの設定（Session メニューの Set Working Directory）を、文字化けするときは fileEncoding を確かめます。Windows の Excel で保存した CSV は fileEncoding = "CP932" で読めることがあります。

\section{metabin で統合する}
\label{sec:05-6}

2 値のアウトカムを統合する関数が metabin です。介入群の事象数、介入群の人数、対照群の事象数、対照群の人数の順に列名を渡し、続けて設定を書きます。この 4 つは、\kkref{sec:04-2} の 2×2 表（表 4.1）では、降圧薬 A 群の行の「脳卒中あり」のマス（ev\_t）と合計のマス（n\_t）、プラセボ群の行の「脳卒中あり」のマス（ev\_c）と合計のマス（n\_c）に当たります。「脳卒中なし」のマスは渡しません。設定のうち、sm = "RR" はリスク比、method = "Inverse" は逆分散法、method.tau = "REML" は τ² の推定法の指定です。

\begin{rcode}
m <- metabin(ev_t, n_t, ev_c, n_c,
             studlab = study, data = bp,
             sm = "RR",            # 効果指標はリスク比
             method = "Inverse",   # 逆分散法
             method.tau = "REML",  # τ² の推定法
             common = TRUE, random = TRUE,
             prediction = TRUE)    # 予測区間も出す
summary(m)
\end{rcode}

出力の後半に、統合した結果が出ます（一部を抜き出し、列の間隔を詰めています）。

\begin{rout}
                         RR           95% CI      z  p-value
Common effect model  0.7481 (0.6724 to 0.8324) -5.33 < 0.0001
Random effects model 0.7281 (0.6137 to 0.8637) -3.64   0.0003
Prediction interval         (0.4600 to 1.1523)

 tau^2 = 0.0336 (0.0000 to 0.2950); tau = 0.1833
 I^2 = 50.8% (0.0% to 76.1%); H = 1.43 (1.00 to 2.05)

     Q d.f. p-value
 18.30    9  0.0318
\end{rout}

最初に見るのは、Random effects model の行です。ランダム効果モデルで統合したリスク比は 0.73（95\% 信頼区間 0.61〜0.86）で、脳卒中の発症がプラセボに比べて平均で 27\% 少ないという推定です。信頼区間が 1 を含まないので、この差が偶然だけで生じたとは考えにくいと判断できます（信頼区間の意味は \kkref{sec:01-1} で説明しました）。固定効果モデル（Common effect model）の 0.75（0.67〜0.83）より区間が広いのは、試験間の真の効果のばらつき（τ² = 0.034）も不確かさに数えるからです。Prediction interval から下の行（tau\^{}2、I\^{}2、H、Q）は \kkref{sec:06-1} で読むので、初めて読むときは飛ばして構いません。

\begin{hattenbox}
method = "MH" とすると、固定効果モデルに Mantel-Haenszel 法を使います。事象の少ない試験が多いときに逆分散法より安定し、RevMan の既定値でもあります。metabin は method を書かないと method = "MH" になります。このとき Mantel-Haenszel 法になるのは固定効果モデルだけで、ランダム効果モデルは逆分散法のままです。この章のコードで method = "Inverse" と書いたのはそのためです。この例の固定効果の値は、逆分散法で 0.7481、Mantel-Haenszel 法で 0.7456 と小数第 3 位から違うだけで、小数第 2 位までは同じ 0.75（0.67〜0.83）です。片方の群で事象が 0 件の試験は、各セルに 0.5 を足して（連続補正）計算されます。両群とも 0 件の試験は効果の向きの情報を持たないので、既定では統合から外れます。

事象がまれなときの方法は、Cochrane ハンドブック（10.4.4）が整理しています。事象の割合が 1\% を下回り、2 群の人数が大きく違わず、効果もあまり大きくないときは、Peto 法の偏りが最も小さく、検出力も高いと報告されています（Bradburn 2007）。Peto 法は、各試験で観察された事象数と、群の差がないと仮定したときに期待される事象数との差から、オッズ比を直接求める方法です。連続補正が要らないので、事象が 0 件の群があっても使えますが、オッズ比にしか使えません（metabin では method = "Peto"）。事象がそれより多いときは、連続補正をしない Mantel-Haenszel 法が勧められます。0.5 を足す連続補正は、推定値を差なしの方向へ偏らせる欠点があります。
\end{hattenbox}


\section{τ² の推定法と信頼区間の作り方}
\label{sec:05-7}

ランダム効果モデルでは、τ² をデータから推定してから重みを決めます。この章では REML という推定法を使いました。試験が少ないときは、区間が狭くなりすぎるのを補う HKSJ 法の区間も並べて比べます。降圧薬 A の 10 試験では、REML の通常の区間が 0.61〜0.86、HKSJ 法の区間が 0.59〜0.89 で、HKSJ 法のほうが広くなりました。

\subsection{\hatten 推定法と区間の作り方の詳細}
\label{sec:05-7-1}

ランダム効果モデルでは、τ² をデータから推定してから重みを決めます。推定法にはいくつかあり、論文の方法の欄に書くべき選択です。古くから使われてきたのが DerSimonian-Laird 法（DL 法、DerSimonian 1986）で、多くのソフトの既定値でしたが、試験数が少ないと τ² を小さく見積もりやすいことが知られています。現在は REML（制限付き最尤法）や Paule-Mandel 法が勧められることが多く、meta パッケージの既定値でこの例でも使った REML は、その一つです。

信頼区間の作り方にも選択肢があります。正規分布を使う標準的な方法は、試験数が少ないと区間が狭くなりすぎ、95\% のはずの区間が真の値を含む割合がそれより低くなります。Hartung-Knapp-Sidik-Jonkman 法（HKSJ 法。Hartung-Knapp 調整とも呼ばれます）は、各試験の推定値が統合値からどれだけ散らばっているかを使って標準誤差を計算し直し、正規分布の代わりに t 分布で区間を作る方法です。t 分布は t 検定で使うものと同じで、試験の数（自由度）が少ないほど区間が広くなります。この方法は、試験数が少ないときの区間の狭さを補います（Hartung 2001、Sidik 2002、IntHout 2014）。

\begin{rcode}
m_hk <- update(m, method.random.ci = "HK")
m_hk
\end{rcode}

\begin{rout}
                         RR           95% CI    z|t  p-value
Random effects model 0.7281 (0.5930 to 0.8938) -3.50   0.0067
\end{rout}

update は、前の結果の設定を一部だけ変えて計算し直す関数です。点推定は変わらず、95\% 信頼区間は 0.59〜0.89 に広がり、p 値は 0.0003 から 0.0067 になりました。HKSJ 法は区間を広げることが多い方法ですが、いつも適切とは限りません。Cochrane ハンドブック（10.10.4）によると、試験が 2 本か 3 本しかないときは区間が広すぎることがあり、τ² が 0 と推定されたときは逆に狭すぎることがあります。後者への対策として、HKSJ 法の標準誤差が標準的な方法より小さくなるときは標準的な方法の値を使う修正も提案されています（IntHout 2014。meta では adhoc.hakn.ci 引数）。どちらの区間を主解析にするかは、データを見る前に計画へ書き、他方は感度分析として並べます。

\section{フォレストプロットを読む}
\label{sec:05-8}

統合の結果は、フォレストプロットという図で示します。forest 関数に metabin の結果を渡すと描けます。

\begin{rcode}
forest(m, sortvar = TE,     # 効果の大きさの順に並べる
       label.left  = "降圧薬 A が有利",
       label.right = "プラセボが有利")
\end{rcode}

Mac で日本語が四角に化けるときは、forest の中に fontfamily = "Hiragino Sans" を足します。図 5.3 は、付録 A のスクリプトで列を減らして保存したものです。

\kkfig{p1_forest.pdf}{図 5.3　降圧薬 A の 10 試験のフォレストプロット（架空のデータ）}

図 5.3 は、上から次の部品でできています。

\begin{description}
\item[試験の行] 1 行が 1 試験です。四角が点推定、左右の横線が 95\% 信頼区間を表します。四角の面積は重みに比例するので、Garcia 2018 や Ivanov 2021 の四角が大きく描かれています。横線が縦の実線（リスク比 1、効果なし）と交わらない試験は、その試験だけで有意な差を示したことになります。
\item[菱形] 統合した結果です。菱形の中心が点推定、左右の端が 95\% 信頼区間です。上が固定効果モデル、下がランダム効果モデルで、ランダム効果モデルの菱形のほうが横に長くなっています。
\item[予測区間の帯] 菱形の下の横帯が予測区間で、読み方は \kkref{sec:06-3} で説明します。菱形よりずっと広く、右端がリスク比 1 を越えています。その下に I²、τ²、異質性の検定の p 値が出ます。
\item[縦の点線] 統合した効果の位置を示す補助線です。横軸は対数目盛（\kkref{sec:05-1-2}）で、右端の列は各試験のリスク比、95\% 信頼区間、ランダム効果モデルでの重みです。
\end{description}

10 試験のうち信頼区間が 1 をまたがないのは 4 試験だけで、残りの 6 試験は単独では有意差なしでした。それでも点推定はほとんどが 1 より左にあり、統合すると菱形は 1 から離れました（\kkref{sec:01-3}）。一方で Chen 2012 のように点推定が 1 より右の試験もあります。

\section{\hatten L'Abbé プロットで両群の発症割合を見る}
\label{sec:05-9}

フォレストプロットはリスク比だけを示すので、各試験で脳卒中がどれくらい起きていたのかは分かりません。L'Abbé（ラベ）プロットは、各試験の 2 つの群の発症割合そのものを、1 試験を 1 点として置く図です。

\begin{rcode}
labbe(m, random = TRUE, common = FALSE)
\end{rcode}

\kkfig[.7\linewidth]{x1_labbe.pdf}{図 5.4　降圧薬 A の 10 試験の L'Abbé プロット。円の大きさは重み、細い実線は両群が等しい線、破線は統合したリスク比（0.73）の線}

横軸はプラセボ群の発症割合、縦軸は降圧薬 A 群の発症割合です。円の中心がその試験の 2 つの割合の組で、円の面積は重みに対応します。左下から右上への細い実線は、両群の発症割合が等しい（リスク比 1）位置です。この線より下にある試験は降圧薬 A 群のほうが発症が少なく、上にある試験は多いことを表します。破線は、統合したリスク比 0.73 に対応する線で、プラセボ群の割合に 0.73 をかけた位置を通ります。

図 5.4 では、ほとんどの試験が実線の下にあり、降圧薬 A が脳卒中を減らす方向にそろっていることが一目で分かります。実線の上にあるのは Chen 2012 だけで、\kkref{sec:01-2} で見た点推定 1.13 に対応します。点が破線のまわりでどれくらい散らばるかを見れば、異質性の大きさの見当もつきます（\kkref{sec:06-1}）。また、試験ごとの対照群の発症割合が、横軸の上でおよそ 5\% から 13\% に広がっていることも、この図ではっきりします。リスク比が同じでも、発症割合が高い試験ほど実線からの距離が大きく、防げる人数は多くなります。この関係は、\kkref{sec:07-9} の結果の要約表で、対照群のリスクを置いて人数に直す理由になっています。

\section{この章のまとめ}
\label{sec:05-10}

効果指標にはリスク比を選び、対数をとって、推定値の分散の逆数を重みにして平均します。主解析は、真の効果が試験ごとに違うと考えるランダム効果モデルとし、τ² は REML で推定します。試験数が少ないときの区間の狭さは、HKSJ 法の区間を並べて確かめます。metabin の出力では、まずランダム効果モデルの行を見ます。どちらを主解析にするかは、データを見る前に計画へ書きます。降圧薬 A の 10 試験では、リスク比は 0.73（95\% 信頼区間 0.61〜0.86）、HKSJ 法の区間は 0.59〜0.89 でした。フォレストプロットでは、試験ごとの結果と菱形、予測区間をあわせて読みます。\kkchref{6}では、試験間のばらつきを調べます。
